\section{Kimi Linear 与 KDA：更细粒度的记忆更新}

本章进入 Kimi Linear 的核心模块 KDA，即 Kimi Delta Attention。节目中提到，KDA 基于 Gated DeltaNet 等早期工作，并把粗粒度 decay 改为更细粒度 decay：过去一个 attention head 下不同维度共享一个衰减率，现在每个维度可以有自己的衰减率。直觉是让有限 hidden state 更充分地表达不同时间尺度的记忆。

\begin{figure}[H]
\centering
\includegraphics[width=0.96\textwidth]{kimi-linear-kda.png}
\caption{Kimi Linear / KDA：Kimi Delta Attention 用更细粒度 decay 改善记忆利用。自制概念图，依据 00:14:39--00:18:56 对谈内容整理。}
\end{figure}

\begin{knowledgebox}{读图：KDA 是旧思想的新组合}
Gated DeltaNet 提供基础，粗粒度 decay 保证效率，细粒度 decay 增强表达，KDA 重新组织记忆更新，最后进入 hybrid 架构。很多新架构不是凭空发明，而是把旧机制在新硬件和新规模下重新组合。
\end{knowledgebox}

\subsection{Scaling Ladder：小规模通关再放大}

本节解释 Kimi 如何降低架构试错风险。线性注意力不是写完公式就能直接上大模型，必须经过逐级规模验证。

Kimi 内部使用 Scaling Ladder：一个模块先在小规模下和 Full Attention 对比，表现足够好才进入更大规模继续 scale。这个过程像通关，避免直接在最大规模上烧算力试错。KDA 的设计也是经过多轮混合方式筛选和规模放大后，逐步成为候选。

\begin{figure}[H]
\centering
\includegraphics[width=0.96\textwidth]{scaling-ladder.png}
\caption{Scaling Ladder：小规模通关后再上更大规模，与 Full Attention 对比。自制概念图，依据 00:18:56--00:20:20 对谈内容整理。}
\end{figure}

\begin{knowledgebox}{读图：架构不是一次性拍脑袋}
小规模筛选降低成本，指标达标后进入更大规模，再与 Full Attention 基线比较。只有成本、性能和 scaling 表现都过关，才可能成为上线候选。
\end{knowledgebox}

\subsection{三比一的混合比例}

接下来讨论 hybrid 架构里的比例问题。保留多少全局层、替换多少线性层，本质上是在表达能力和效率之间找一个可扩展的工程点。

访谈提到，Kimi Linear 每三层 KDA 插入一层全注意力层，三比一比例逐渐成为一种共识。这个比例的本质是用 Full Attention 维护全局表达能力，用 KDA 降低大部分层的推理成本。它不是数学定理，而是工程实验中形成的经验折中。

\begin{knowledgebox}{KDA 机制消化}
\begin{tabularx}{\textwidth}{p{0.22\textwidth}p{0.34\textwidth}X}
\toprule
机制 & 作用 & 为什么重要 \\
\midrule
Delta Rule & 用递推方式更新记忆状态 & 让 attention 具有类似 RNN 的状态更新能力。 \\
Gating & 控制信息保留和遗忘 & 让模型选择哪些历史信息应该留下。 \\
Decay & 记忆衰减率 & 不同信息可以有不同时间尺度。 \\
Fine-grained decay & 每个维度有独立衰减 & 更充分利用有限 hidden state。 \\
Hybrid 全局层 & 定期插入 Full Attention & 兜底全局表达和长程交互。 \\
\bottomrule
\end{tabularx}
\end{knowledgebox}

\begin{warningbox}{KDA 不是“把 Softmax 删掉”这么简单}
如果只去掉 Softmax，模型可能表达力不足。KDA 的价值在于重新设计状态更新、门控和衰减，让线性层尽量保留有用历史信息；Full Attention 层则负责补全全局交互。
\end{warningbox}

\subsection{本章小结}

Kimi Linear 的关键不是“全部换成线性注意力”，而是在 hybrid 架构中用 KDA 提供高效记忆更新，再用少量全局层保证下限。Scaling Ladder 则是让这个设计能逐步走向大规模训练的机制。

